\documentclass[a4paper,12pt]{article}
\usepackage[utf8]{inputenc}
\usepackage[T1]{fontenc}
\usepackage[portuguese]{babel}
\usepackage{geometry}
\usepackage[dvipsnames]{xcolor}
\usepackage{soul}
\usepackage{hyperref}
\usepackage{microtype}
\geometry{margin=2.5cm}

% \hl usa soul que permite quebra de linha dentro do highlight
\sethlcolor{yellow}

\title{Primeira Liga}
\author{Trabalho Prático 1 --- Scripting no PLN}
\date{2025/26}

\begin{document}

\maketitle

\begin{abstract}
As seguintes frases foram selecionadas através de um modelo de linguagem
baseado em \textit{n-grams} (trigramas) com \textit{scoring} por
log-probabilidade média e suavização de Laplace:
\begin{itemize}
  \item Afinal, veio parar ao Leixões, que não sente nenhuma honra com o facto.
  \item Na época 2011-12 foi introduzido um novo troféu maior e mais pesado que o original e que era atribuído ao clube que desde então conseguisse ganhar três campeonatos consecutivos ou cinco campeonatos intercalados.
  \item Na época 2014-15 regressou-se ao modelo de dezoito equipas, motivada pelas pretensões de vários clubes de menor dimensão bem como pela integração do Boavista Futebol Clube, devido à prescrição do procedimento disciplinar ocorrido em 2008, devido ao processo Apito Final.
\end{itemize}
\end{abstract}

\section{Texto da Fonte}

A \hl{Primeira Liga} (\hl{Liga Portugal Betclic} por razões de patrocínio) é o principal escalão do sistema de ligas de futebol de \hl{Portugal}. Criada na época 1934-35 pela \hl{Federação Portuguesa de Futebol}, é organizada pela \hl{Liga Portuguesa de Futebol Profissional} desde a temporada 1995-96. É disputada por dezoito clubes, num sistema de promoção e despromoção com a \hl{Segunda Liga}.
A partir da época 2023-24 a \hl{Primeira Liga} tem o nome comercial de \hl{Liga Portugal Betclic} devido a um acordo de patrocínio entre a casa de apostas e a \hl{Liga Portugal}, fechando assim um ciclo de dois anos em que a bwin foi o principal patrocinador da prova.
No final da época 2023-24 a \hl{Liga Portuguesa} ocupava o 7.º lugar no ranking da \hl{UEFA}, o que significa que na época seguinte apenas o campeão tem acesso direto à \hl{Liga dos Campeões} e o 2.º classificado à 3.ª pré-eliminatória da competição. Os 3.º e 4.º classificados têm acesso, respetivamente, às 2.ª pré-eliminatórias da \hl{Liga Europa} e da \hl{Liga Conferência}.
Os clubes classificados em 17.º e 18.º lugares são despromovidos à \hl{Segunda Liga}, por troca com os 1.º e 2.º classificados desta prova que são assim promovidos a primodivisionários. A equipa que terminar em 16.º lugar disputará um play-off de despromoção/promoção a duas mãos com o 3.º lugar da \hl{Segunda Liga}.
Adicionalmente, as equipas da \hl{Primeira Liga} participam na \hl{Taça de Portugal} e na \hl{Taça da Liga}, entrando na 3.ª e 2.ª eliminatórias destas competições, respetivamente.
Durante as 91 edições disputadas até ao momento, participaram na \hl{Primeira Liga} um total de 73 clubes, dos quais somente cinco se sagraram campeões nacionais. O maior vencedor da história da \hl{Primeira Liga} é o \hl{Benfica}, com 38 campeonatos nacionais conquistados. O atual campeão nacional é o \hl{Sporting}, após conquistar na época 2024-25 o seu 21.º título.
Em 1921, após a derrota da \hl{Seleção Nacional} na sua estreia frente à \hl{Espanha}, surgiu a necessidade de se alterar o sistema do futebol português, constituído por campeonatos regionais (\hl{Porto} e \hl{Lisboa}, com algumas competições irregulares na \hl{Madeira}). Nasceu assim uma prova regular com os vencedores das provas distritais chamada \hl{Campeonato de Portugal}, prova que em 1938 passaria depois a designar-se \hl{Taça de Portugal}. Na sua primeira edição, na época 1921-22, teve apenas dois clubes, \hl{Sporting} e \hl{FC Porto} (a representar \hl{Lisboa} e \hl{Porto}, respetivamente). O \hl{FC Porto} venceu numa finalíssima a primeira edição. Em 1934 começou então verdadeiramente o \hl{Campeonato Nacional da Primeira Divisão}, com oito equipas, catorze jornadas a duas voltas e a somar pontos, e em que o \hl{FC Porto} foi o primeiro vencedor. Na altura foi chamada de \hl{Liga Experimental}, tendo em conta que era a primeira vez que se organizava. Foi assim, a partir da época 1934-35, que os campeões nacionais passaram a ser designados a partir do \hl{Campeonato da Liga da Primeira Divisão} (época 1934-35) e que até hoje já teve cinco vencedores. A competição anterior, o chamado \hl{Campeonato de Portugal}, era uma prova por eliminatórias, incluindo clubes da \hl{Segunda Divisão} cujos vencedores eram definidos numa final (no entanto, os títulos dos \hl{Campeonatos} de \hl{Portugal} não contam como títulos da \hl{Taça de Portugal}, nem de títulos do \hl{Campeonato da Primeira Divisão} de acordo com o que ficou definido no \hl{Relatório de Atividades} da FPF de 1938).
O surgimento do \hl{Campeonato da Primeira Divisão} teve muito que ver com uma nova derrota sofrida pela seleção nacional em \hl{Madrid} por 9-0, no apuramento para o \hl{Mundial de 1934}, em que várias vozes questionaram a competitividade do modelo do \hl{Campeonato de Portugal}, nomeadamente o número reduzido de jogos disputados por cada equipa e o valor dos competidores em prova. \hl{Ricardo Ornelas} escreveu no jornal \hl{Os Sports} que se deveria realizar uma prova em poule, à semelhança do que acontecia na principais potências futebolísticas da \hl{Europa}. No sentido de aumentar a competitividade do futebol português, a FPF encarregou \hl{Plácido de Souza}, \hl{Ribeiro dos Reis}, \hl{Cândido de Oliveira} e \hl{Virgílio da Fonseca} de elaborarem o projeto de uma nova competição em poule. No entanto, por causa da situação económica do país, a FPF tinha dúvidas sobre a viabilidade económica da prova, devido às deslocações a que os participantes estariam sujeitos, bem como sobre o acolhimento que teria junto do público. Na época 1934-35 foi criado o \hl{Campeonato da Liga da Primeira Divisão}. Após o sucesso da competição, em 1938 a FPF decidiu o seguinte:
"Por virtude da reforma a que se procedeu no \hl{Estatuto e Regulamentos da Federação} os \hl{Campeonatos das Ligas} e de \hl{Portugal} passaram a designar-se, respectivamente, \hl{Campeonatos Nacionais} e \hl{Taça de Portugal}".
Ao vencedor do \hl{Campeonato da Liga da Primeira Divisão} (competição organizada a título experimental mas cujos títulos são considerados oficiais) seria atribuído o título de campeão nacional.
Participaram nesta primeira edição oito clubes na \hl{Primeira Divisão} (quatro de \hl{Lisboa}, dois do \hl{Porto}, um de \hl{Coimbra} e um de \hl{Setúbal} - os campeonatos regionais mais competitivos da época).
O sucesso da prova foi imediato, não só económico mas sobretudo desportivo, com a sucessão de jogos disputados pelas melhores equipas, o que levou a que popularmente a prova relegasse para um plano secundário o \hl{Campeonato de Portugal}. O jornalista \hl{Ricardo Ornelas} por mais de uma ocasião no jornal \hl{Os Sports} defendeu que o vencedor da \hl{Liga} é que deveria ser considerado campeão nacional. Mais tarde tal viria a acontecer por parte da FPF.
Num congresso realizado em agosto de 1938 dá-se uma remodelação dos regulamentos das provas da FPF, em que ficou estabelecido:
"acabar com os \hl{Campeonatos das Ligas} e substituir o \hl{Campeonato de Portugal} das jornadas em sucessiva eliminações, por um campeonato de maior rigor e regularidade, pelo sistema de "poule" em duas voltas"
Na prática traduziu-se apenas em renomear o "\hl{Campeonato da Liga da Primeira Divisão}" para "\hl{Campeonato Nacional da Primeira Divisão}" (sendo a principal categoria muitas vezes abreviada para "\hl{Primeira Divisão}") e renomearam o "\hl{Campeonato de Portugal}" para "\hl{Taça de Portugal}", de acordo com o \hl{Relatório de Atividades} 1938 da FPF. A designação manteve-se até 1999, tendo nessa altura o nome sido alterado para "\hl{Primeira Liga}".
O \hl{Futebol Clube do Porto} foi o primeiro vencedor do campeonato, numa altura em que se disputava entre oito equipas. \hl{Manuel Soeiro}, jogador do \hl{Sporting Clube de Portugal} foi o primeiro melhor marcador do campeonato, com catorze golos em catorze jogos. O \hl{Sporting}, que ficou a dois pontos do campeão nessa época, só venceu a liga na época 1940-41, já na época da \hl{Primeira Divisão}.
Em 1935-36, foi a vez do \hl{Benfica} se sagrar campeão, por três vezes consecutivas. O \hl{Clube de Futebol} Os \hl{Belenenses} foi o quarto campeão diferente da liga, vencida na época 1945-46. No século seguinte, foi a vez do \hl{Boavista Futebol Clube} inscrever-se na lista de campeões de \hl{Portugal}. Desta vez, o clube portuense venceu a liga na época 2000-01.
"\hl{Os Três Grandes}" é uma expressão que tradicionalmente designa os três principais clubes de futebol \hl{em Portugal}: \hl{Benfica}, \hl{Porto} e \hl{Sporting}. Estes são os clubes com mais títulos de campeão nacional e, igualmente, com mais segundos e terceiros lugares. Juntos "\hl{Os Três Grandes}" detêm 89 dos 91 títulos de campeão disputados (1934-35 até 2024-25): o \hl{Benfica} tem 38 títulos, o \hl{Porto} tem 30 títulos e o \hl{Sporting} tem 21 títulos. Nas 91 épocas completas já disputadas na \hl{Primeira Liga}, em 54 temporadas o pódio foi exclusivamente ocupado pelos \hl{Três Grandes}.
O campeonato iniciou-se na época 1934-35 e confrontou apenas oito equipas na \hl{Primeira Divisão}: os quatro primeiros classificados do campeonato regional de \hl{Lisboa}, os dois melhores do \hl{Porto}, o campeão de \hl{Setúbal} e o campeão de \hl{Coimbra} (os quatro campeonatos regionais mais competitivos) enquanto as restantes equipas dos regionais eram apuradas para a \hl{II Divisão}. O início da época 1939-40 ficou marcada pela polémica, devido a uma batalha administrativa entre o \hl{FC Porto} e o \hl{Académico} do \hl{Porto} relativamente a um jogo do \hl{Campeonato Regional do Porto}. A \hl{Federação Portuguesa de Futebol} arranjou uma solução para satisfazer os dois clubes, alargando o campeonato para 10 equipas.
Um jogo do \hl{Campeonato Regional da AF Porto} entre o \hl{FC Porto} e o \hl{Académico Futebol Clube} acabou sendo interrompido pelo árbitro após um anormal número de expulsões e lesões, sobretudo do lado do \hl{FC Porto}, atribuindo a vitória ao \hl{Académico}. No entanto a decisão acabou sendo contestada pelo \hl{FC Porto}, dado que os regulamentos da altura não previam a interrupção do jogo por número mínimo de participantes e a \hl{AF Porto} deliberou a repetição do jogo, que resultou em vitória do \hl{FC Porto}.
O \hl{Campeonato} terminaria com \hl{FC Porto} em primeiro, seguido de \hl{Leixões SC} e \hl{Académico}. No entanto, este último recorreu da decisão da \hl{AF Porto} para a FPF. Dada a polémica instalada, a FPF decidiu pelo alargamento da \hl{Primeira Divisão} para dez clubes, abrindo-se uma vaga extra para a \hl{AF Porto} e outra para a \hl{AF Setúbal}, decisão que teria o voto contra do \hl{FC Porto}, segundo os dirigentes do \hl{Académico}, para impedir a participação deste no campeonato, dada a animosidade:
...como se sabe o [FC] \hl{Porto} votou contra a inclusão de mais um grupo tripeiro só para nos prejudicar, o que sendo uma deslealdade, é um tanto anti-bairrista.
Para além disso, a FPF anulou também o jogo de repetição entre \hl{FC Porto} e \hl{Académico}, o que relegou o \hl{FC Porto} para a 3ª posição do campeonato regional, e atribuiu automaticamente o título regional ao \hl{Leixões SC}, que no entanto repudiou publicamente a situação:
O \hl{Leixões} repudia a benesse. O meu clube não aceita título que não ganhou! O \hl{Leixões} não quer ser campeão por favor. Não lhe assenta bem um título usurpado a outrem. \hl{Acho} que foi infeliz a decisão da FPF! O \hl{FC Porto} não merecia semelhante castigo, apenas para ser beneficiado um terceiro. Afinal, veio parar ao \hl{Leixões}, que não sente nenhuma honra com o facto.
Na época seguinte, a prova voltaria a ser disputada por oito equipas. Na época 1941-42 foi decidido que o campeonato seria alargado de oito para dez equipas para admitir os campeões da \hl{AF Braga} e \hl{AF Algarve} (até esta época apenas os dois primeiros classificados dos campeonatos regionais das AFs do \hl{Porto}, \hl{Coimbra}, \hl{Lisboa} e \hl{Setúbal} eram admitidos). O \hl{FC Porto} acabou o campeonato regional em terceiro lugar, o que não dava acesso à \hl{Primeira Divisão}. Contudo, um segundo alargamento (de dez para doze equipas) na mesma época foi decidido, o que permitiu ao clube participar na \hl{Primeira Divisão}. Este número de clubes ir-se-ia manter até à época 1945-46, altura em que admitiu doze equipas (entraram os campeões da \hl{AF Évora} e da \hl{AF Aveiro}).
Na época 1946-47, dá-se uma reformulação dos quadros competitivos, acabando-se com a qualificação a partir dos campeonatos regionais, passando a existir uma lógica de continuidade entre edições, e um sistema de promoções e descidas entre divisões. A \hl{Primeira Divisão} foi alargada para catorze equipas, enquanto a \hl{II Divisão} foi reformulada, e criada uma \hl{III Divisão}.
O número de equipas na \hl{Primeira Divisão} manteve-se durante vinte e cinco épocas, até que na época 1971-72 passou a dezasseis equipas para na época 1987-88 passar a admitir vinte, assim se mantendo por duas épocas. Na época 1989-90 assume o formato das dezoito equipas, com uma exceção na temporada seguinte (vinte), mantendo-se assim até à época 2005-06, sendo que na época 2006-07 houve uma redução para dezasseis equipas.
Na época 2014-15 regressou-se ao modelo de dezoito equipas, motivada pelas pretensões de vários clubes de menor dimensão bem como pela integração do \hl{Boavista Futebol Clube}, devido à prescrição do procedimento disciplinar ocorrido em 2008, devido ao processo \hl{Apito Final}. Optou-se portanto pelo arquivamento, sem qualquer juízo sobre a existência ou não da infração que pendia sobre o \hl{Boavista}. Desta maneira impôs-se a sua reintrodução na \hl{Primeira Liga}.
Em consequência da pandemia de \hl{COVID-19}, após considerar inicialmente a realização de jogos à porta fechada, a \hl{Liga Portuguesa de Futebol Profissional} decidiu a 12 de março de 2020 pela suspensão total dos jogos da \hl{Primeira Liga} na época 2019-20 por tempo indeterminado. A competição foi retomada a partir de 3 de \hl{Junho} de 2020, com os jogos disputados à porta fechada.
O troféu de campeão nacional é entregue anualmente pela FPF, também a \hl{Liga} entrega em cada época um troféu ao vencedor da \hl{Primeira Liga}.
Na época 2011-12 foi introduzido um novo troféu maior e mais pesado que o original e que era atribuído ao clube que desde então conseguisse ganhar três campeonatos consecutivos ou cinco campeonatos intercalados. Este troféu foi apenas entregue ao SL \hl{Benfica} pelos campeonatos ganhos nas épocas 2013-14, 2014-15 e 2015-16 e a partir da época 2016-17 deixou de ser entregue.
O acesso às competições de clubes da \hl{UEFA} é feito tendo por base a posição da \hl{Primeira Liga} no ranking da \hl{UEFA}. Presentemente, fruto do 7.º lugar no ranking, \hl{Portugal} tem uma vaga direta na fase de grupos da \hl{Liga dos Campeões}, para o campeão nacional, enquanto que o segundo classificado terá acesso à 3.ª pré-eliminatória. O vencedor da \hl{Taça de Portugal} terá acesso direto à fase de grupos da \hl{Liga Europa}. Já os 3.º e 4.º lugares darão acesso, respetivamente, às 2.ª pré-eliminatórias da \hl{Liga Europa} e da \hl{Liga Conferência}. Contudo, se o vencedor da \hl{Taça de Portugal} tiver conseguido a qualificação para a \hl{Liga dos Campeões} através do \hl{Campeonato}, o 3.º classificado é apurado para a fase de grupos da \hl{Liga Europa} e o 4.º e 5.º classificados para as 2.ª pré-eliminatórias da \hl{Liga Europa} e da \hl{Liga Conferência}.

\bibliographystyle{plain}
\begin{thebibliography}{1}
  \bibitem{wikipedia_primeira_liga}
  Primeira Liga.
  Disponível em: \url{https://pt.wikipedia.org/wiki/Primeira\_Liga}.
  Acedido em março de 2026.
\end{thebibliography}

\end{document}
